\subsection{Folgen und Grenzwerte in $\mathbb R^n$}
$\lim_{x\rightarrow a}f(x)=y$, wenn für alle Folgen $\{x^k\} \subset \Omega \setminus \{a\}$ mit $\lim_{k\rightarrow \infty}x^k=a$ gilt:
\begin{equation*}
    \lim_{k\rightarrow \infty}f(x^k)=a
\end{equation*}
Findet man eine Folge, wo dies nicht der Fall ist, so kann $a$ nicht der Grenzwert sein

\subsection{Stetigkeit}
Skalarfeld $f:\Omega \subset \mathbb R^n \rightarrow \mathbb R$ stetig an $a \in \Omega$, falls
\begin{equation*}
    \lim_{x\rightarrow a}f(x)=f(a)
\end{equation*}

\subsection{Wichtige Sätze}
\begin{itemize}
    \item $f:\Omega \rightarrow \mathbb R$ stetig und $\Omega$ kompakte Menge $\Longrightarrow f$ nimmt auf $\Omega$ ihr Minimum und Maximum an
\end{itemize}